% ============================================================
\section{Metodi scelti e documentazione}
% ============================================================

\subsection{BrokerImpl}

Il Javadoc di classe di \texttt{BrokerImpl} chiarisce la separazione di responsabilità
della classe: gestisce il comportamento a livello di oggetto, mentre delega allo
\texttt{StoreManager} (fornito al momento dell'inizializzazione) tutta l'interazione con
lo store.

La documentazione usata per determinare i metodi target è stata principalmente quella della classe target, delle
interfacce che implementa direttamente e quelle che queste ultime estendono
transitivamente, escludendo i tipi di libreria standard per rimanere concentrati sulle interfacce del progetto.
L'analisi dei metodi è stata condotta a best effort. Il focus è stato dedicato maggiormente a \texttt{Broker} e \texttt{StoreContext} in quanto
ipotizzate come più rilevanti.

Gli attributi di \texttt{BrokerImpl} sono stati consultati per identificare i parametri di stato del SUT.
Alcuni metodi variano anche lo stato di oggetti passati in input, quindi anche questi sono stati analizzati.
Anche il manuale OpenJPA e la documentazione online sono state utilizzate, tuttavia descrivono \texttt{BrokerImpl}
senza aggiungere informazioni sul comportamento dei metodi.

La scelta dei metodi da testare è stata guidata bilanciando la presenza di \textbf{parametri di input},
la \textbf{significatività} del metodo e la disponibilità di \textbf{documentazione} (non è sempre
possibile disporre di un oracolo completo). Per gestire la grande mole di metodi, una prima scrematura è stata effettuata
ordinando in modo decrescente i metodi per linee di javadoc. I metodi selezionati, con l'interfaccia o la classe che li dichiara, sono riportati in
Tabella~\ref{tab:brokerimpl-metodi} (Appendice).

\subsubsection{Documentazione}
\label{sec:brokerimpl-analisi-doc}

\begin{itemize}[itemsep=12pt, parsep=7pt]
    \item\texttt{find(Object oid, boolean validate, FindCallbacks call)}

    \textit{Documentato esplicitamente:} restituisce l'istanza della classe persistente
    associata all'oid dato. Il Javadoc distingue tre casi:
        \begin{itemize}
        
        \item \texttt{validate=true} e oggetto assente $\to$ null;
        \item \texttt{validate=true} e presente $\to$ non-null;
        \item\texttt{validate=false} $\to$ mai null, restituendo l'istanza in cache se presente,
        altrimenti tentando di crearne una hollow, altrimenti \texttt{ObjectNotFoundException}.
    
        \end{itemize}
    Se validate è true, e viene convalidata la presenza dell'oggetto nello store, vengono caricati i fetch group fields.

    \textit{Dedotto / documentato altrove:} il parametro \texttt{validate} ``convalida'' la presenza dell'oggetto nello store.
    \texttt{StoreContext} espone anche un secondo overload di \texttt{find}, con Javadoc che documenta un comportamento di default legato a
    \texttt{flags}, controllabile tramite le costanti: queste ultime
    non hanno però una propria documentazione. Il parametro \texttt{exclude} ha invece un
    valore documentato, \texttt{StoreContext.EXCLUDE\_ALL}.
    \texttt{BrokerImpl} espone inoltre un terzo overload, \texttt{protected}, il cui Javadoc
    (``Internal finder.'') non aggiunge informazioni utili oltre al nome. Il parametro
    \texttt{call} è di tipo
    \texttt{FindCallbacks}: il suo Javadoc di interfaccia documenta \texttt{processArgument}
    (può rifiutare l'argomento con un'eccezione, sostituirlo con un altro id, o ignorarlo
    restituendo \texttt{null}) e \texttt{processReturn} (riceve l'id e uno state manager, il
    suo \texttt{@return} dichiara il valore ``da restituire'').

    \textit{Riferimenti:} \texttt{StoreContext.java} (\texttt{find} a 3 e 5 parametri,
    \texttt{EXCLUDE\_ALL}), \texttt{BrokerImpl.java} (\texttt{find} \texttt{protected}),
    \texttt{FindCallbacks.java}.

    \item\texttt{lock(Object pc, int level, int timeout, OpCallbacks call)}

    \textit{Documentato esplicitamente:} assicura che l'istanza indicata venga bloccata al
    livello richiesto. Il Javadoc documenta solo \texttt{timeout} come millisecondi di
    attesa prima di rinunciare (\texttt{-1} per attesa illimitata); non contiene alcun
    riferimento esplicito a dove siano definiti i livelli ammessi né a come verificarne
    l'esito.

    \textit{Dedotto / documentato altrove:}
    \texttt{Broker} estende \texttt{LockLevels}, che dichiara i livelli ammessi; \texttt{StoreContext}
    dichiara inoltre \texttt{getLockLevel}, da cui si deduce che la postcondizione di
    \texttt{lock} sia verificabile con \texttt{getLockLevel(pc) == level}. Il parametro
    \texttt{call} è di tipo \texttt{OpCallbacks}: il suo unico metodo,
    \texttt{processArgument}, può rifiutare l'argomento con un'eccezione o restituire un
    codice azione, senza che
    l'effetto sia documentato specificamente per \texttt{lock}. \texttt{Broker} espone anche
    un secondo overload, \texttt{lock(Object pc, OpCallbacks call)}: utile per sapere che il
    livello di lock ha anche una fonte alternativa (quello corrente della
    \texttt{FetchConfiguration} del broker), non usata nell'overload qui scelto.
    \texttt{OpenJPAStateManager} (tipo del terzo parametro di
    \texttt{OpCallbacks.processArgument}, quindi collegato strutturalmente a \texttt{call})
    espone \texttt{getPCState()}; \texttt{PCState} dichiara con Javadoc esplicito il
    significato di ciascuno stato (\texttt{TRANSIENT}: ``unmanaged instance''; \texttt{PNEW}:
    ``Persistent-New''; \texttt{PCLEAN}/\texttt{PDIRTY}: ``Persistent-Clean''/``Persistent-Dirty'';
    \texttt{HOLLOW}: ``exists in data store''). \texttt{PCState} dichiara in totale 23 costanti:
    delle restanti, nessuna ha un appiglio testuale rilevante per questo metodo (dettaglio in
    Sezione~\ref{sec:cp}). \texttt{StoreManager} (raggiungibile via
    \texttt{StoreContext.getStoreManager()}) conferma il significato di \texttt{HOLLOW} (metodo
    \texttt{load}) e che le istanze \texttt{PNEW} non hanno ancora una riga fisica nello store
    prima del flush (metodo \texttt{flush}). Questi stati definiscono le classi di \texttt{pc}.
    \texttt{LockManager} (raggiungibile via \texttt{StoreContext.getLockManager()}) documenta
    \texttt{@param level}: ``one of the lock constants defined in \texttt{LockLevels}, or a
    custom level'' — distingue esplicitamente costanti dichiarate (incluso \texttt{LOCK\_NONE},
    elencato nello stesso gruppo testuale di \texttt{LOCK\_READ}/\texttt{LOCK\_WRITE}, e quelle
    di \texttt{MixedLockLevels}, che estende \texttt{LockLevels}) da valori non noti arbitrari;
    questi due gruppi definiscono le classi di \texttt{level}.

    \textit{Riferimenti:} \texttt{Broker.java} (\texttt{lock} a 4 e 2 parametri),
    \texttt{LockLevels.java}, \texttt{MixedLockLevels.java}, \texttt{StoreContext.java}
    (\texttt{getLockLevel}, \texttt{getStoreManager}, \texttt{getLockManager}),
    \texttt{OpCallbacks.java}, \texttt{OpenJPAStateManager.java} (\texttt{getPCState}),
    \texttt{PCState.java}, \texttt{StoreManager.java} (\texttt{load}, \texttt{flush}),
    \texttt{LockManager.java} (\texttt{lock}).

    \item\texttt{detachAll(Collection objs, OpCallbacks call)}

    \textit{Documentato esplicitamente:} stacca dal broker le istanze indicate, producendo
    copie non gestite. Garantisce il valore di ritorno e specifica che vengono attraversati
    solo i campi della \texttt{FetchConfiguration} corrente; per staccare un grafo di
    oggetti collegati, le relazioni verso altre istanze persistenti devono trovarsi nel
    \texttt{default-fetch-group} o nella \texttt{FetchConfiguration} personalizzata corrente.

    \textit{Dedotto / documentato altrove:} tramite collegamento con il Javadoc di
    \texttt{Broker.isDetached(Object obj)} si può affermare che ogni istanza restituita
    soddisfa \texttt{isDetached(e)=true}. \texttt{Broker} espone altri due overload di
    \texttt{detachAll}: rivelano un comportamento di flush automatico (sempre eseguito, o reso
    opzionale) non presente nell'overload con collezione qui scelto.
    Il tipo dichiarato (\texttt{Collection}) permette di partizionare anche
    quanti elementi contiene la collezione, oltre a cosa può essere un singolo elemento. Il
    parametro \texttt{call} è di tipo \texttt{OpCallbacks} (vedi \texttt{lock}). Come per
    \texttt{lock}/\texttt{persist}, un singolo elemento della collezione è classificato tramite
    \texttt{PCState}: \{non gestito\} (\texttt{TRANSIENT}) contro \{gestito\} (ogni altro stato,
    non ulteriormente distinto poiché \texttt{detachAll} non documenta un effetto diverso tra
    sotto-stati gestiti). Aggiunta anche una classe con elementi misti (un gestito e un non
    gestito nella stessa collezione), per testare l'atomicità non documentata dell'operazione.

    \textit{Riferimenti:} \texttt{Broker.java} (\texttt{detachAll} nei suoi 3 overload,
    \texttt{isDetached(Object)}), \texttt{OpCallbacks.java}, \texttt{OpenJPAStateManager.java}
    (\texttt{getPCState}), \texttt{PCState.java}.

    \item\texttt{attachAll(Collection objs, boolean copyNew, OpCallbacks call)}

    \textit{Documentato esplicitamente:} importa nel broker le istanze indicate. Distingue
    due comportamenti in base allo stato di ogni istanza: se precedentemente staccata da
    questo o da un altro broker, le modifiche vengono unite (merge) nell'istanza persistente
    esistente; se nuova, viene persistita come tale; \texttt{copyNew} controlla
    esplicitamente questo secondo caso.

    \textit{Dedotto / documentato altrove:} come per \texttt{detachAll}, il tipo
    \texttt{Collection} permette di partizionare anche la cardinalità, oltre al contenuto di un
    singolo elemento. Il parametro \texttt{call} è di tipo \texttt{OpCallbacks} (vedi
    \texttt{lock}). A differenza di \texttt{detachAll}, qui \texttt{TRANSIENT} non è il caso non
    valido: coincide col ramo documentato ``nuovo''. Il caso non valido è invece
    \texttt{PCState.PDELETED} (``Persistent-Deleted''): un'istanza già marcata per cancellazione,
    concettualmente incompatibile con un attach/merge, non coperta da nessuno dei due rami
    documentati. Aggiunta anche una classe con elementi misti (un valido e un non valido),
    stesso motivo di \texttt{detachAll}.

    \textit{Riferimenti:} \texttt{Broker.java} (\texttt{attachAll}), \texttt{OpCallbacks.java},
    \texttt{OpenJPAStateManager.java} (\texttt{getPCState}), \texttt{PCState.java}.

    \item\texttt{newInstance(Class cls)}

    \textit{Documentato esplicitamente:} crea una nuova istanza di tipo \texttt{cls}.
    Documenta due esiti: se \texttt{cls} è un'interfaccia o una classe astratta con metodi
    conformi alla convenzione JavaBeans, ne genera un'implementazione concreta dai metadati;
    se \texttt{cls} è già un tipo managed concreto, ne restituisce direttamente un'istanza;
    altrimenti lancia \texttt{IllegalArgumentException}.

    \textit{Riferimenti:} \texttt{Broker.java} (\texttt{newInstance}).

    \item\texttt{isDetached(Object obj, boolean find)}

    \textit{Documentato esplicitamente:} verifica se l'oggetto fornito è staccato. Se
    \texttt{find=true}, come ultima risorsa verifica anche l'esistenza dell'oggetto nel
    database, considerandolo staccato se presente. Il Javadoc non descrive esplicitamente
    il comportamento per \texttt{find=false}.

    \textit{Dedotto / documentato altrove:} \texttt{Broker} dichiara un overload a
    un solo parametro, \texttt{isDetached(Object obj)}, il cui Javadoc definisce
    esplicitamente cosa significhi ``detached'': un'istanza ``che può
    essere riagganciata a un Broker tramite una chiamata a \texttt{Broker\#attach}''.
    Come per \texttt{lock}/\texttt{detachAll}/\texttt{persist}, \texttt{obj} non gestito è
    definito tramite \texttt{PCState.TRANSIENT}; ``detached'' invece non è uno dei 23 stati di
    \texttt{PCState} (verificato), è definito solo dal Javadoc di \texttt{isDetached} stesso.

    \textit{Riferimenti:} \texttt{BrokerImpl.java} (\texttt{isDetached} a 2 parametri),
    \texttt{Broker.java} (\texttt{isDetached} a 1 parametro),
    \texttt{OpenJPAStateManager.java} (\texttt{getPCState}), \texttt{PCState.java}.

    \item\texttt{persist(Object pc, Object id, OpCallbacks call)}

    \textit{Documentato esplicitamente:} rende persistente l'istanza data; il \texttt{@return}
    documenta che restituisce lo state manager assegnato alla nuova istanza. Documenta
    \texttt{id} come opzionale (\texttt{null} assegna un identificativo di default) e
    chiarisce che, a differenza di altre operazioni di persist, questo metodo non effettua
    cascata immediata sui campi marcati \texttt{CASCADE\_IMMEDIATE}.

    \textit{Dedotto / documentato altrove:} \texttt{Broker} espone anche
    \texttt{persist(Object, OpCallbacks)} (Javadoc di una riga, nessuna postcondizione utile) e
    \texttt{persistAll}. \texttt{ValueMetaData}
    documenta un terzo valore di cascata oltre a \texttt{CASCADE\_NONE}/
    \texttt{CASCADE\_IMMEDIATE} (\texttt{CASCADE\_AUTO}), e chiarisce che
    \texttt{CASCADE\_NONE} genera un errore se una relazione transiente è ancora presente al
    momento del flush: possibile scenario di eccezione.
    Il parametro \texttt{call} è di tipo
    \texttt{OpCallbacks} (vedi \texttt{lock}). Come per \texttt{lock}, \texttt{pc} è definito
    tramite \texttt{PCState}/\texttt{getPCState()}: qui però il caso documentato dal Javadoc è
    quello opposto, un'istanza \texttt{TRANSIENT} (``unmanaged instance'') che diventa
    persistente; il caso di un \texttt{pc} già gestito (\texttt{PNEW}/\texttt{PCLEAN}/
    \texttt{PDIRTY}/\texttt{HOLLOW}) non ha alcun riscontro testuale.

    \textit{Riferimenti:} \texttt{Broker.java} (\texttt{persist} nei suoi 3 overload),
    \texttt{ValueMetaData.java} (\texttt{CASCADE\_*}), \texttt{OpCallbacks.java},
    \texttt{OpenJPAStateManager.java} (\texttt{getPCState}), \texttt{PCState.java}.
\end{itemize}

\subsection{LoginDialog}

% analisi da svolgere